\section*{Capítulo II, Problema VII}

\textbf{Problema:}

7. Encuentra la ecuación del plano perpendicular al vector dado $N$ y que pasa por el punto dado $P$.

\begin{itemize}
    \item[(a)] $N = (1, -1, 3)$, $P = (4, 2, -1)$
    \item[(b)] $N = (-3, -2, 4)$, $P = (2, \pi, -5)$
    \item[(c)] $N = (-1, 0, 5)$, $P = (2, 3, 7)$
\end{itemize}

\textbf{Solución:}

Si un plano es perpendicular al vector $N = (a,b,c)$ y pasa por $P = (x_0,y_0,z_0)$, entonces $N$ es un vector normal del plano y su ecuación es
\[
    a(x-x_0) + b(y-y_0) + c(z-z_0) = 0.
\]

\begin{itemize}
    \item[(a)] Con $N = (1,-1,3)$ y $P = (4,2,-1)$:
    \[
        (x-4) - (y-2) + 3(z+1) = 0.
    \]
    Simplificando,
    \[
        x - y + 3z + 1 = 0.
    \]

    \item[(b)] Con $N = (-3,-2,4)$ y $P = (2,\pi,-5)$:
    \[
        -3(x-2) - 2(y-\pi) + 4(z+5) = 0.
    \]
    Simplificando,
    \[
        -3x - 2y + 4z + 26 + 2\pi = 0.
    \]

    \item[(c)] Con $N = (-1,0,5)$ y $P = (2,3,7)$:
    \[
        -(x-2) + 5(z-7) = 0.
    \]
    Simplificando,
    \[
        -x + 5z - 33 = 0.
    \]
\end{itemize}